% Research draft produced for the accompanying request, 20 September 2026.
\documentclass[11pt,a4paper]{article}
\usepackage[T1]{fontenc}
\usepackage[utf8]{inputenc}
\usepackage{lmodern}
\usepackage[margin=2.45cm]{geometry}
\usepackage{amsmath,amssymb,amsthm,mathtools}
\usepackage{microtype,booktabs,array,longtable}
\usepackage{xcolor}
\usepackage{enumitem}
\usepackage{fancyhdr,needspace}
\usepackage[colorlinks=true,linkcolor=blue!45!black,citecolor=blue!45!black,
  urlcolor=blue!45!black]{hyperref}
\hypersetup{pdftitle={Rectangular Schur Polynomials and Cigler's Conjecture 8},
 pdfauthor={AI-assisted research draft},
 pdfsubject={An even-shift Hankel determinant identity and its consequences}}
\newtheorem{theorem}{Theorem}[section]
\newtheorem{lemma}[theorem]{Lemma}
\newtheorem{proposition}[theorem]{Proposition}
\newtheorem{corollary}[theorem]{Corollary}
\theoremstyle{definition}
\newtheorem{definition}[theorem]{Definition}
\theoremstyle{remark}
\newtheorem{remark}[theorem]{Remark}
\newcommand{\CT}{\operatorname{CT}_z}
\newcommand{\Q}{\mathbb Q}
\newcommand{\Z}{\mathbb Z}
\newcommand{\C}{\mathbb C}
\newcommand{\N}{\mathbb N}
\newcommand{\SL}{\operatorname{SL}}
\newcommand{\GL}{\operatorname{GL}}
\newcommand{\Sym}{\operatorname{Sym}}
\newcommand{\diag}{\operatorname{diag}}
\newcommand{\len}{\ell}
\newcommand{\Rkn}{R_{k,n}}
\newcommand{\Xk}{X_k}
\DeclareMathOperator{\sgn}{sgn}
\numberwithin{equation}{section}
\setlist{nosep}
\emergencystretch=2em
\allowdisplaybreaks[2]
\pagestyle{fancy}
\fancyhf{}
\fancyhead[L]{\small Rectangular Schur polynomials and Hankel determinants}
\fancyhead[R]{\small Research draft}
\fancyfoot[C]{\thepage}
\setlength{\headheight}{14pt}
\title{\textbf{Rectangular Schur Polynomials\protect\\
and Cigler's Conjecture 8}\\[0.4em]
\large A Laurent--Gram proof of the even-shift Hankel identity}
\author{AI-assisted research draft}
\date{20 September 2026}

\begin{document}
\maketitle
\begin{abstract}
Let
\(b_m(t)=\sum_{r\ge0}\binom{\lfloor m/2\rfloor}{r}
\binom{\lceil m/2\rceil}{r}t^r\).
We prove, for all integers \(k,n\ge0\), the polynomial identity
\[
 \det\bigl(b_{2k+i+j}(t)\bigr)_{0\le i,j<n}
 =t^{\lfloor n^2/4\rfloor}
 s_{(n^k)}(\underbrace{1,\ldots,1}_{k},
                \underbrace{t,\ldots,t}_{k}).
\]
The proof realizes the Hankel matrix as a Laurent-polynomial Gram matrix,
changes to a consecutive monomial basis, and identifies the resulting
Toeplitz determinant by dual Jacobi--Trudi. A positive sum over partitions
inside a rectangle proves full coefficient support and palindromicity;
restriction of a Schur module to \(\SL_2\) proves unimodality. A confluent
alternant supplies explicit numerator polynomials of exactly the degrees and
with exactly the reciprocity required in Cigler's Conjecture~8
(arXiv:2111.14492v3). Thus every assertion of that conjecture is obtained,
not only its positivity assertion. Further consequences include stable edge
coefficients, a dimension product, an exact generic recurrence, and
fixed-parameter asymptotics. All identities and boundary cases are proved
symbolically; a standard-library Python verifier supplies additional exact
finite checks.
\end{abstract}

\medskip
\noindent\textbf{Status.}
This is a complete mathematical proof draft relative to the classical Schur
identities and Schur-module character theorem stated below. It has not been
independently refereed or checked in a proof assistant. A targeted public
literature search did not locate an earlier resolution of the selected
conjecture; this is not an exhaustive priority certificate.

\clearpage
\tableofcontents
\clearpage

\section{The selected problem and its provenance}
\label{sec:problem}

\subsection{Why this is related to, but not a duplicate of, the manifest}
The supplied manifest \cite{manifest} contains an entry entitled
\emph{Product formulas for ballot-polynomial Hankel determinants}. Its stated
scope is Conjectures~13--15 in Section~5 of Cigler's paper \cite{cigler}.
The present article instead treats \emph{Conjecture~8 in Section~4}, concerning
the polynomial extension \(b_m(t)\) of the middle binomial coefficients.
The determinant family, normalization, and numbered conjecture are all
specified below. None of the manifest's reports is used as a theorem in the
proof. In particular, the manifest's descriptions of claimed results are not
silently treated as independently verified mathematics.

The original source is the third version of arXiv:2111.14492, dated
30 December 2021. The arXiv record consulted on 20 September 2026 still
identified that version as current. Conjecture~8 appears on printed
pages~16--17, with its numerator formula and reciprocity numbered (56)
and (57). A targeted search for the paper title and for combinations of
``Cigler'', ``Conjecture 8'', ``Schur'', and ``proof'' found no prior
resolution. An absence of search results cannot establish novelty.
The package's \texttt{sources.md} records the search and its limits.

\subsection{Notation and the precise target}
For integers \(m\ge0\), define
\begin{equation}
 b_m(t)=\sum_{r=0}^{\lfloor m/2\rfloor}
 \binom{\lfloor m/2\rfloor}{r}
 \binom{\lceil m/2\rceil}{r}t^r.
 \label{eq:moment}
\end{equation}
Thus
\[
 b_0=b_1=1,\qquad b_2=1+t,\qquad b_3=1+2t,\qquad
 b_4=1+4t+t^2.
\]
At \(t=1\), Vandermonde's convolution gives
\(b_m(1)=\binom{m}{\lfloor m/2\rfloor}\), explaining the name.
For a shift \(s\ge0\) and an order \(n\ge0\), put
\begin{equation}
 D_s(n,t)=\det\bigl(b_{s+i+j}(t)\bigr)_{0\le i,j<n},
 \qquad D_s(0,t)=1.
 \label{eq:hankel}
\end{equation}
The empty determinant convention is important throughout. We will prove
\begin{equation}
 D_0(n,t)=t^{\lfloor n^2/4\rfloor}.
 \label{eq:unshifted}
\end{equation}
For even shifts we write
\begin{equation}
 R_{k,n}(t)=\frac{D_{2k}(n,t)}{D_0(n,t)}
 =\frac{D_{2k}(n,t)}{t^{\lfloor n^2/4\rfloor}}.
 \label{eq:normalization}
\end{equation}
At this point the quotient is interpreted in \(\Q(t)\); its polynomiality
is a conclusion, not an assumption. In Cigler's notation,
\(R_{k,n}(t)=d_{2k}(n,t)=r_{2k}(n,t)\).

For \(k\ge1\), Conjecture~8 asks that \(R_{k,n}\) have positive coefficients
at every degree from \(0\) through \(kn\), be palindromic and unimodal, and
admit the expression
\begin{equation}
 R_{k,n}(t)=
 \frac{\displaystyle\sum_{j=0}^{k}(-1)^j
 c_{2k,j}(n,t)t^{j(n+j)}}{(1-t)^{k^2}},
 \label{eq:target-numerator}
\end{equation}
where
\begin{align}
 \deg_t c_{2k,j}(n,t)&=2j(k-j),\label{eq:target-degree}\\
 c_{2k,k-j}(n,t)&=t^{2j(k-j)}c_{2k,j}(n,t^{-1}).
 \label{eq:target-reciprocity}
\end{align}
Here \emph{unimodal} means weakly increasing up to the center and weakly
decreasing afterwards; strict unimodality is neither required nor true in
general. We prove all of these assertions, including the stated degrees.

\section{Schur notation and the classical ingredients}
\label{sec:schur}
A partition is a weakly decreasing finite list of nonnegative integers,
with trailing zeros ignored. Its size is \(|\lambda|=\sum_i\lambda_i\),
and its conjugate is denoted \(\lambda'\). The notation \((n^k)\) means the
rectangle with \(k\) rows, each of length \(n\). If \(k=0\) or \(n=0\),
this is the empty partition.

For a list \(X=(x_1,\ldots,x_m)\), let \(e_r(X)\) and \(h_r(X)\) be the
elementary and complete homogeneous symmetric polynomials:
\begin{equation}
 \sum_{r\ge0}e_r(X)u^r=\prod_{a=1}^m(1+x_au),\qquad
 \sum_{r\ge0}h_r(X)u^r=\prod_{a=1}^m(1-x_au)^{-1}.
 \label{eq:eh}
\end{equation}
Both are zero for negative subscripts and equal one at subscript zero;
\(e_r(X)=0\) also when \(r>m\).

A semistandard Young tableau has weakly increasing rows and strictly
increasing columns. The Schur polynomial \(s_\lambda(X)\) is the sum of the
weight monomials of all such tableaux of shape \(\lambda\), with entries
from \(\{1,\ldots,m\}\). The following classical descriptions of the same
polynomial will be used \cite{fulton,macdonald}:
\begin{align}
 s_\lambda(X)
 &=\det\bigl(h_{\lambda_i-i+j}(X)\bigr)_{1\le i,j\le \len(\lambda)},
 \label{eq:jt}\\
 s_\lambda(X)
 &=\det\bigl(e_{\lambda'_i-i+j}(X)\bigr)_{1\le i,j\le\lambda_1},
 \label{eq:dual-jt}\\
 s_\lambda(X)
 &=\frac{\det(x_i^{\lambda_{m+1-a}+a-1})_{1\le i,a\le m}}
         {\det(x_i^{a-1})_{1\le i,a\le m}}.
 \label{eq:alternant}
\end{align}
The last formula is first read for distinct variables and then extended as
a polynomial identity. The increasing order of the exponents in
\eqref{eq:alternant} fixes the signs used later.

We also use one standard representation-theoretic fact. If \(V\) is a
complex vector space of dimension \(m\), the Schur module
\(S_\lambda(V)\) is a homogeneous polynomial representation of
\(\GL(V)\) of degree \(|\lambda|\), and the trace of
\(\diag(x_1,\ldots,x_m)\) on it is \(s_\lambda(X)\).
The tableau basis has exactly the corresponding weight monomials. This
character theorem is explicitly stated in Fulton's treatment
\cite{fulton}. It is an imported classical result; it is not claimed as a
new construction here. The needed elementary \(\SL_2\) argument is proved
in Section~\ref{sec:unimodal}.

For our repeated-variable specialization, write
\begin{equation}
 X_k=(\underbrace{1,\ldots,1}_{k},
        \underbrace{t,\ldots,t}_{k}).
 \label{eq:alphabet}
\end{equation}
For example,
\begin{align}
 e_r(X_k)&=\sum_a\binom{k}{a}\binom{k}{r-a}t^a,\label{eq:e-explicit}\\
 h_r(X_k)&=\sum_{a=0}^{r}\binom{k+a-1}{a}
                   \binom{k+r-a-1}{r-a}t^a \quad(k\ge1,\ r\ge0).
 \label{eq:h-explicit}
\end{align}
We use the usual convention that a binomial coefficient with an invalid
lower index is zero.

\section{The central determinant identity}
\label{sec:main}
\begin{theorem}[Even-shift Hankel--Schur identity]
\label{thm:main}
For all integers \(k,n\ge0\),
\begin{equation}
 \boxed{\displaystyle
 D_{2k}(n,t)=t^{\lfloor n^2/4\rfloor}s_{(n^k)}(X_k).}
 \label{eq:main}
\end{equation}
The identity holds in \(\Z[t]\). In particular,
\(D_0(n,t)=t^{\lfloor n^2/4\rfloor}\) and
\(R_{k,n}(t)=s_{(n^k)}(X_k)\).
\end{theorem}

The determinant on the left has entries depending on \(i+j\). The proof
turns it into a determinant depending on \(i-j\). The entire normalizing
power of \(t\) comes from the change of basis, which is why both parities
of the matrix order must be tracked explicitly.

\subsection{A constant-term model for the moments}
Work temporarily over \(K=\Q(t)\) and in the Laurent-polynomial ring
\(K[z,z^{-1}]\). Define
\begin{equation}
 A(z)=(1+z)(1+t/z)=1+t+z+t/z.
 \label{eq:A}
\end{equation}
Let \(\CT\) extract the coefficient of \(z^0\). Substitution
\(z\mapsto t/z\) defines a \(K\)-linear ring involution \(f\mapsto f^*\).
It fixes \(A\), and
\begin{equation}
 \CT f^*=\CT f.
 \label{eq:ct-star}
\end{equation}
Indeed, the coefficient of \(z^0\) is unchanged when each monomial
\(z^r\) is replaced by \(t^rz^{-r}\).

\begin{lemma}
\label{lem:moments}
For \(a\ge0\),
\begin{equation}
 b_{2a}(t)=\CT A^a,\qquad
 b_{2a+1}(t)=\CT(1+z)A^a
            =\CT(1+t/z)A^a.
 \label{eq:ct-moments}
\end{equation}
\end{lemma}
\begin{proof}
Expand \(A^a=(1+z)^a(1+t/z)^a\). Its constant term is
\(\sum_r\binom ar^2t^r\). Similarly, the constant term of
\((1+z)^{a+1}(1+t/z)^a\) is
\(\sum_r\binom{a+1}{r}\binom ar t^r\). These are the even and odd cases
of \eqref{eq:moment}. The last equality follows from
\eqref{eq:ct-star}.
\end{proof}

\subsection{The Hankel matrix as a Gram matrix}
Define a \emph{linear} map \(\Phi:K[x]\to K[z,z^{-1}]\) on monomials by
\begin{equation}
 \Phi(x^{2a})=A^a,\qquad
 \Phi(x^{2a+1})=(1+z)A^a.
 \label{eq:Phi}
\end{equation}
This map is not an algebra homomorphism; no multiplicativity of \(\Phi\)
is used. For fixed \(k\ge0\), define the bilinear form
\begin{equation}
 B_k(f,g)=\CT A^k f(z)g(t/z)
 \quad\text{on }K[z,z^{-1}].
 \label{eq:form}
\end{equation}
It is symmetric because applying the involution inside the constant term
interchanges \(f\) and \(g\).

Put \(v_i=\Phi(x^i)\). Then
\begin{equation}
 B_k(v_i,v_j)=b_{2k+i+j}(t).
 \label{eq:gram}
\end{equation}
For completeness, if \(i=2a,j=2b\), the product inside the constant term
is \(A^{k+a+b}\). If \(i=2a+1,j=2b+1\), it is
\((1+z)(1+t/z)A^{k+a+b}=A^{k+a+b+1}\). The two mixed-parity cases give
\((1+z)A^{k+a+b}\) and \((1+t/z)A^{k+a+b}\), respectively. Lemma~\ref{lem:moments}
proves \eqref{eq:gram} in all four cases.

Consequently, \(D_{2k}(n,t)\) is the Gram determinant of
\(v_0,\ldots,v_{n-1}\) for the form \(B_k\).

\subsection{A triangular Laurent basis change}
Consider the ordered monomials
\begin{equation}
 w_0=1,\ w_1=z,\ w_2=z^{-1},\ w_3=z^2,\ w_4=z^{-2},\ldots;
 \qquad w_{2a}=z^{-a},\quad w_{2a+1}=z^{a+1}.
 \label{eq:laurent-basis}
\end{equation}
Each \(v_i\) belongs to the span of \(w_0,\ldots,w_i\). The coefficient
of the newly appearing monomial is
\begin{equation}
 [w_{2a}]v_{2a}=t^a,\qquad [w_{2a+1}]v_{2a+1}=1.
 \label{eq:diagonal}
\end{equation}
For the even case, \(A^a\) has lowest power \(z^{-a}\), with coefficient
\(t^a\), and highest power \(z^a\). For the odd case,
\((1+z)A^a\) has highest power \(z^{a+1}\), with coefficient one, and
lowest power \(z^{-a}\). This proves both triangularity and
\eqref{eq:diagonal}.

Assume \(n\ge1\), and set \(q=\lfloor(n-1)/2\rfloor\). The determinant
of the transition matrix from the first \(n\) monomials \(w_i\) to the
first \(n\) vectors \(v_i\) is therefore
\[
 \prod_{a=0}^{q}t^a=t^{q(q+1)/2}.
\]
A bilinear Gram determinant changes by the square of this determinant:
\begin{equation}
 D_{2k}(n,t)=t^{q(q+1)}
 \det\bigl(B_k(w_i,w_j)\bigr)_{0\le i,j<n}.
 \label{eq:gram-change}
\end{equation}

Reorder the monomials increasingly by exponent in both rows and columns.
The two permutation signs multiply to one. The resulting exponent set is
\begin{equation}
 \begin{cases}
 \{-m,-m+1,\ldots,m\},&n=2m+1,\\
 \{-(m-1),-(m-2),\ldots,m\},&n=2m.
 \end{cases}
 \label{eq:exponents}
\end{equation}
For any two such exponents \(a,b\),
\begin{equation}
 B_k(z^a,z^b)=t^b\CT A^kz^{a-b}.
 \label{eq:monomial-gram}
\end{equation}
Factor \(t^b\) from the column indexed by \(b\). The sum of the exponents
in \eqref{eq:exponents} is zero for \(n=2m+1\) and \(m\) for \(n=2m\).
Combining with \eqref{eq:gram-change}, the total power extracted is
\begin{equation}
 \begin{cases}
 m(m+1),&n=2m+1,\\
 (m-1)m+m=m^2,&n=2m,
 \end{cases}
 =\lfloor n^2/4\rfloor.
 \label{eq:power-audit}
\end{equation}
Since the exponent set is consecutive, the remaining determinant is
\begin{equation}
 D_{2k}(n,t)=t^{\lfloor n^2/4\rfloor}
 \det\bigl(\CT A^kz^{i-j}\bigr)_{0\le i,j<n}.
 \label{eq:toeplitz}
\end{equation}
This identity has been obtained in \(\Q(t)\), where all column factors are
legitimate, including those with negative exponent.

\subsection{Identification of the Toeplitz determinant}
We have
\begin{equation}
 z^kA(z)^k=(1+z)^k(z+t)^k
          =\prod_{x\in X_k}(z+x).
 \label{eq:elementary-symbol}
\end{equation}
The coefficient of \(z^{2k-r}\) in this product is \(e_r(X_k)\). Hence
\begin{equation}
 \CT A^kz^{i-j}=e_{k+i-j}(X_k).
 \label{eq:entry-e}
\end{equation}
The conjugate of the rectangle \((n^k)\) is \((k^n)\). After transposing
the determinant, dual Jacobi--Trudi \eqref{eq:dual-jt} gives
\begin{equation}
 \det\bigl(e_{k+i-j}(X_k)\bigr)_{0\le i,j<n}=s_{(n^k)}(X_k).
 \label{eq:toeplitz-schur}
\end{equation}
Equations~\eqref{eq:toeplitz} and \eqref{eq:toeplitz-schur} prove
Theorem~\ref{thm:main} for \(n\ge1\). For \(n=0\), both determinants and
the empty Schur polynomial are one. For \(k=0\), the matrix in
\eqref{eq:toeplitz-schur} is the identity matrix, so
\eqref{eq:unshifted} follows as well. Both sides of \eqref{eq:main} belong
to \(\Z[t]\); equality in \(\Q(t)\) is therefore the claimed polynomial
identity. In particular it remains valid at \(t=0\), without dividing a
specialized zero by zero. This completes the proof.

\section{An explicit positive partition formula}
\label{sec:positive}
For a partition \(\lambda\) with at most \(k\) parts, pad it with zeros
and define
\begin{equation}
 d_k(\lambda)=\prod_{1\le a<b\le k}
 \frac{\lambda_a-\lambda_b+b-a}{b-a}.
 \label{eq:weyl-dimension}
\end{equation}
This is \(s_\lambda(1^k)\), the number of semistandard tableaux of shape
\(\lambda\) on \(k\) letters, and is a positive integer. The product
formula also follows directly from the alternant by letting all variables
tend to one: normalized derivatives give a Vandermonde in the exponents
\(\lambda_{k+1-a}+a-1\), divided by the Vandermonde in \(a-1\).

\begin{theorem}[Positive sum of squared tableau counts]
\label{thm:positive}
For \(k,n\ge0\),
\begin{equation}
 \boxed{\displaystyle
 R_{k,n}(t)=\sum_{\lambda\subseteq(n^k)}d_k(\lambda)^2t^{|\lambda|}.}
 \label{eq:positive}
\end{equation}
For \(k=0\), the only summand is the empty partition, of dimension one.
\end{theorem}
\begin{proof}
Take a tableau of the rectangle \((n^k)\) on the alphabet
\(\{1,\ldots,2k\}\). The boxes whose entries are at most \(k\) form a
Young diagram \(\lambda\subseteq(n^k)\). The entries in those boxes have
\(d_k(\lambda)\) possibilities. The other boxes form the skew shape
\((n^k)/\lambda\), and each contributes a factor \(t\).

Define the rectangular complement
\begin{equation}
 \lambda^c=(n-\lambda_k,n-\lambda_{k-1},\ldots,n-\lambda_1).
 \label{eq:complement}
\end{equation}
Rotate the skew tableau through \(180\) degrees, first subtract \(k\)
from its entries, and replace each resulting letter \(a\) by \(k+1-a\).
Rotation reverses both row and column directions; reversal of the alphabet
restores weak increase along rows and strict increase down columns. This is
a bijection with tableaux of the ordinary shape \(\lambda^c\) on \(k\)
letters. Therefore splitting the alphabet gives
\[
 s_{(n^k)}(1^k,t^k)
 =\sum_{\lambda\subseteq(n^k)}
 d_k(\lambda)d_k(\lambda^c)t^{kn-|\lambda|}.
\]
The product \eqref{eq:weyl-dimension}, with pairs of indices reversed,
shows \(d_k(\lambda^c)=d_k(\lambda)\). Replacing \(\lambda\) by
\(\lambda^c\) in the sum and using Theorem~\ref{thm:main} proves
\eqref{eq:positive}.
\end{proof}

\begin{corollary}[Polynomiality, full support, and symmetry]
\label{cor:support}
For \(k\ge1,n\ge0\), every coefficient of \(R_{k,n}\) from degree zero
through degree \(kn\) is a positive integer. Its degree is \(kn\), its
constant and leading coefficients are one, and
\begin{equation}
 R_{k,n}(t)=t^{kn}R_{k,n}(t^{-1}).
 \label{eq:palindrome}
\end{equation}
\end{corollary}
\begin{proof}
For every \(0\le r\le kn\), there exists a partition of size \(r\) inside
the rectangle: fill as many whole rows of length \(n\) as possible and
then at most one partial row. The case \(n=0\) is immediate. Each
corresponding dimension in \eqref{eq:positive} is positive. Only the empty
partition contributes at degree zero, and only the full rectangle
contributes at degree \(kn\); both have dimension one on \(k\) letters.
Finally \(\lambda\mapsto\lambda^c\) preserves the squared dimension and
changes the size from \(r\) to \(kn-r\).
\end{proof}

This argument explains the nonnegativity without extracting coefficients
from an alternating determinant. It does \emph{not} by itself prove
unimodality: a positive, symmetric sum of monomials need not be unimodal.
The next section supplies the missing mechanism.

\section{Unimodality by a balanced \texorpdfstring{$\SL_2$}{SL2} restriction}
\label{sec:unimodal}
\begin{lemma}[Weight strings]
\label{lem:sl2}
Suppose a finite-dimensional complex algebraic representation of
\(\SL_2\) has torus character
\[
 u^{-N}P(u^2),\qquad P(t)=\sum_{r=0}^{N}a_rt^r.
\]
Then there exist nonnegative integers \(\mu_0,\ldots,
\mu_{\lfloor N/2\rfloor}\) such that
\begin{equation}
 P(t)=\sum_{a=0}^{\lfloor N/2\rfloor}
 \mu_a t^a(1+t+\cdots+t^{N-2a}).
 \label{eq:strings}
\end{equation}
In particular the coefficient sequence, including any zero end coefficients,
is symmetric and unimodal.
\end{lemma}
\begin{proof}
For clarity, recall why finite-dimensional \(\SL_2\) modules have the
required strings. Restriction to the compact group \(\mathrm{SU}_2\)
admits an invariant positive-definite Hermitian form, obtained by averaging
any such form. Orthogonal complements of invariant subspaces are invariant;
complexifying the Lie algebra gives the same complete reducibility for
\(\mathfrak{sl}_2\), and hence for the connected group \(\SL_2\).

With generators \(E,F,H\) satisfying
\([H,E]=2E\), \([H,F]=-2F\), and \([E,F]=H\), choose a highest-weight
vector \(v\) in an irreducible summand, so that \(Ev=0\) and \(Hv=mv\).
The commutator relations imply, by induction,
\[
 EF^rv=r(m-r+1)F^{r-1}v.
\]
If \(F^qv\ne0\) and \(F^{q+1}v=0\), this formula gives \(m=q\).
The span of \(v,Fv,\ldots,F^mv\) is invariant, so irreducibility makes
it the whole summand. Its weights are \(m,m-2,\ldots,-m\), each once.
Thus the summand has character
\(u^{-m}+u^{-m+2}+\cdots+u^m\), as does \(\Sym^m\C^2\).

In the character assumed in the lemma, all weights lie between \(-N\)
and \(N\), and all have parity \(N\). Since a direct sum has nonnegative
weight multiplicities, every occurring \(m\) has the form \(N-2a\), with
\(0\le a\le\lfloor N/2\rfloor\). Multiplying the corresponding character
by \(u^N\) and substituting \(t=u^2\) gives the summand in
\eqref{eq:strings}. All strings have the same center, \(N/2\). Their sum
is consequently symmetric and weakly increases towards that center.
\end{proof}

\begin{theorem}[Unimodality]
\label{thm:unimodal}
For all \(k\ge1,n\ge0\), the polynomial \(R_{k,n}(t)\) is unimodal.
More precisely, with \(N=kn\), it admits \eqref{eq:strings} with
\(P=R_{k,n}\), and
\begin{equation}
 \mu_a=[t^a]R_{k,n}(t)-[t^{a-1}]R_{k,n}(t)\ge0
 \quad(0\le a\le\lfloor N/2\rfloor),
 \label{eq:multiplicities}
\end{equation}
where the coefficient of \(t^{-1}\) is zero.
\end{theorem}
\begin{proof}
Identify \(\C^{2k}\) with \(\C^2\otimes\C^k\). Let \(\SL_2\) act on
the first factor and trivially on the second, then restrict the Schur module
\(S_{(n^k)}(\C^{2k})\) to this subgroup. The torus element
\(\diag(u^{-1},u)\) has \(k\) eigenvalues \(u^{-1}\) and \(k\)
eigenvalues \(u\) on the underlying vector space. Homogeneity of the Schur
polynomial and Theorem~\ref{thm:main} give the restricted character
\[
 s_{(n^k)}((u^{-1})^k,u^k)=u^{-kn}s_{(n^k)}(1^k,(u^2)^k)
                       =u^{-kn}R_{k,n}(u^2).
\]
Apply Lemma~\ref{lem:sl2}. Up to the center, the coefficient of \(t^a\)
in \eqref{eq:strings} is \(\mu_0+\cdots+\mu_a\), proving
\eqref{eq:multiplicities}.
\end{proof}

\begin{remark}
The restriction argument applies to \(s_\lambda(1^k,t^k)\) for any
partition \(\lambda\) with at most \(2k\) parts, if the coefficient list
is extended through degree \(|\lambda|\) by zeros where necessary.
The rectangle is needed for the particular Hankel identity and for the
nonzero endpoints, not for the general weight-string mechanism.
We do not claim this classical representation-theoretic mechanism as new.
\end{remark}

\section{Explicit numerator polynomials and reciprocity}
\label{sec:numerator}
This section proves the remaining assertions
\eqref{eq:target-numerator}--\eqref{eq:target-reciprocity}. All signs and
normalizing factorials are included, because the repeated-variable limit is
otherwise an easy source of indexing errors.

Fix \(k\ge1\) and \(n\ge0\). Index positions by
\(I=\{0,1,\ldots,2k-1\}\), and put
\begin{equation}
 \epsilon_a=\begin{cases}a,&0\le a<k,\\ n+a,&k\le a<2k,\end{cases}
 \qquad F_k=\prod_{r=0}^{k-1}r!.
 \label{eq:epsilon}
\end{equation}
These \(2k\) exponents are strictly increasing, including when \(n=0\).
For a subset \(S\subseteq I\), let \(\epsilon_S\) denote the list in
increasing index order, and write
\[
 \Delta(\epsilon_S)=\prod_{\substack{a,b\in S\\a<b}}
 (\epsilon_b-\epsilon_a).
\]
For a \(k\)-element subset \(S\), set
\begin{align}
 j(S)&=|S\cap\{k,k+1,\ldots,2k-1\}|,\label{eq:jS}\\
 \delta(S)&=\sum_{a\in S}a-\binom{k}{2}-j(S)^2,\label{eq:delta}\\
 W_S(n)&=\frac{\Delta(\epsilon_S)\Delta(\epsilon_{S^c})}{F_k^2}.
 \label{eq:weight}
\end{align}
Finally define
\begin{equation}
 C_{k,j}(n,t)=
 \sum_{\substack{S\subseteq I,\ |S|=k\\j(S)=j}}
 (-1)^{\delta(S)}W_S(n)t^{\delta(S)}
 \qquad(0\le j\le k).
 \label{eq:blocks}
\end{equation}
This is an explicit finite sum with \(\binom{2k}{k}\) terms in total.

\begin{lemma}[Bounds and integrality]
\label{lem:blocks}
If \(j(S)=j\), then
\begin{equation}
 0\le\delta(S)\le 2j(k-j).
 \label{eq:delta-bounds}
\end{equation}
Each endpoint is achieved by a unique such \(S\). Moreover, \(W_S(n)\)
is a positive integer for every integer \(n\ge0\). As a polynomial in
\(n\), it belongs to \(\Q[n]\) and has exact degree \(2j(k-j)\).
\end{lemma}
\begin{proof}
The smallest sum of indices is obtained by taking the smallest \(k-j\)
indices from the lower half of \(I\) and the smallest \(j\) from the upper
half. Their sum is \(\binom{k}{2}+j^2\). The largest sum, obtained by
taking the largest indices in the same two halves, is larger by
\(2j(k-j)\). These extremizers are unique, proving the bounds.

For increasing nonnegative integers \(q_1<\cdots<q_k\),
\begin{equation}
 \det\left(\binom{q_b}{r}\right)_{
      \substack{0\le r<k\\1\le b\le k}}
 =\frac{\prod_{a<b}(q_b-q_a)}{F_k}.
 \label{eq:binomial-vandermonde}
\end{equation}
Indeed, the polynomial \(\binom{x}{r}\) has degree \(r\) and leading
coefficient \(1/r!\). Changing from these polynomials to powers of \(x\)
in the determinant proves \eqref{eq:binomial-vandermonde}. Its left side
is an integer, and its right side is strictly positive. Applying it to
\(\epsilon_S\) and \(\epsilon_{S^c}\) proves the integrality and
positivity of \(W_S(n)\).

Only differences between one lower-half exponent and one upper-half
exponent depend on \(n\). There are \(j(k-j)\) such differences in each
of the two Vandermonde products. Each has the form \(n+c\), with \(c>0\).
The claimed degree is therefore exact and its leading coefficient is positive.
\end{proof}

\begin{theorem}[The conjectured numerator, explicitly]
\label{thm:numerator}
For all \(k\ge1,n\ge0\),
\begin{equation}
 \boxed{\displaystyle
 (1-t)^{k^2}R_{k,n}(t)
 =\sum_{j=0}^{k}(-1)^j t^{j(n+j)}C_{k,j}(n,t).}
 \label{eq:numerator}
\end{equation}
For each such \(n\), the polynomials \(C_{k,j}(n,t)\) have integer
coefficients and exact degree \(2j(k-j)\) in \(t\). They satisfy
\begin{equation}
 C_{k,k-j}(n,t)=t^{2j(k-j)}C_{k,j}(n,t^{-1}),
 \qquad C_{k,0}=C_{k,k}=1.
 \label{eq:C-reciprocity}
\end{equation}
As a bivariate expression \(C_{k,j}\in\Q[n,t]\), its degree in \(n\) is
also exactly \(2j(k-j)\). Thus the requested numerator polynomials are
\(c_{2k,j}=C_{k,j}\).
\end{theorem}

\subsection{Confluence of the alternant}
Use the alternant \eqref{eq:alternant} for the partition \((n^k)\) padded
to length \(2k\). Its increasing exponents are precisely
\(\epsilon_0,\ldots,\epsilon_{2k-1}\).
Let the first \(k\) variables tend to one and the last \(k\) tend to \(t\).
Divide numerator and denominator by the within-group Vandermonde products
before taking the limit. The normalized derivative of order \(r\) of
\(x^{\epsilon_a}\) at \(x=y\) is
\(\binom{\epsilon_a}{r}y^{\epsilon_a-r}\).
The remaining denominator is the product of the \(k^2\) cross-group
differences, namely \((t-1)^{k^2}\). Therefore
\begin{equation}
 R_{k,n}(t)=\frac{\det M}{(t-1)^{k^2}},
 \label{eq:confluent}
\end{equation}
where \(M\) is the \(2k\)-by-\(2k\) matrix with columns indexed by \(a\in I\):
\begin{align}
 M_{r,a}&=\binom{\epsilon_a}{r},&&0\le r<k,\label{eq:M-top}\\
 M_{k+r,a}&=\binom{\epsilon_a}{r}t^{\epsilon_a-r},&&0\le r<k.
 \label{eq:M-bottom}
\end{align}
An entry with \(r>\epsilon_a\) is zero. This computation is performed in
\(\Q(t)\), so there is no need to assign a meaning to an intermediate
negative power at \(t=0\).

The normalized-derivative rule used here follows by Taylor expansion:
when \(k\) row variables coalesce at \(y\), the lowest alternating term
in their differences is the Vandermonde times the determinant of derivative
rows of orders \(0,\ldots,k-1\), each divided by its factorial. Thus the
calculation does not assume any analytic convergence of infinite series.

\subsection{Laplace expansion and its sign}
Expand \(\det M\) along the bottom \(k\) rows, selecting columns \(S\).
The top minor is \(\Delta(\epsilon_{S^c})/F_k\). Factoring powers of
\(t\) from the columns and rows of the bottom minor gives
\begin{equation}
 \det M_{\text{bottom},S}
 =t^{\sum_{a\in S}\epsilon_a-\binom{k}{2}}
   \frac{\Delta(\epsilon_S)}{F_k}.
 \label{eq:bottom-minor}
\end{equation}
By the definitions of \(j\) and \(\delta\), the exponent is
\begin{equation}
 \sum_{a\in S}\epsilon_a-\binom{k}{2}
 =jn+j^2+\delta(S).
 \label{eq:exponent-split}
\end{equation}
Using zero-based row and column indices, the Laplace sign is
\((-1)^{k(3k-1)/2+\sum_{a\in S}a}\). Changing the denominator in
\eqref{eq:confluent} from \((t-1)^{k^2}\) to \((1-t)^{k^2}\) adds the
factor \((-1)^{k^2}\). The resulting sign simplifies as follows:
\begin{align}
 (-1)^{k^2+k(3k-1)/2+\sum_{a\in S}a}
 &=(-1)^{\binom{k}{2}+\sum_{a\in S}a}\notag\\
 &=(-1)^{j^2+\delta(S)}=(-1)^{j+\delta(S)}.
 \label{eq:sign}
\end{align}
Grouping subsets according to \(j(S)\) proves \eqref{eq:numerator}.

\subsection{Exact degrees and complementary subsets}
Lemma~\ref{lem:blocks} shows that \(C_{k,j}\) is an integer polynomial
in \(t\), with degree at most \(d=2j(k-j)\). The unique subset at
\(\delta=d\) contributes the nonzero coefficient
\((-1)^dW_S(n)=W_S(n)>0\). Hence the degree is exactly \(d\), for every
nonnegative integer \(n\). The unique subset at \(\delta=0\) likewise
shows that the coefficient of \(n^d\) in \(C_{k,j}\), as a polynomial
in \(t\), has nonzero constant term. This proves the exact degree in \(n\).

Complementing the selected subset gives
\begin{equation}
 j(S^c)=k-j(S),\qquad
 \delta(S^c)=2j(k-j)-\delta(S),\qquad W_{S^c}(n)=W_S(n).
 \label{eq:complement-subset}
\end{equation}
Since \(2j(k-j)\) is even, the signs of complementary terms agree.
Equation~\eqref{eq:C-reciprocity} follows immediately. If \(j=0\) or
\(j=k\), there is only one subset, both Vandermonde products are those
of consecutive integers, and their quotient by \(F_k^2\) is one.
This completes the proof of Theorem~\ref{thm:numerator}.

\begin{remark}[Integer-valued is not the same as integral coefficients in \(n\)]
The formula gives \(C_{k,j}(n,t)\in\Z[t]\) at every integer \(n\ge0\),
but we only assert \(C_{k,j}\in\Q[n,t]\) before specializing \(n\).
The Vandermonde quotients are integer-valued polynomials. No claim that
their ordinary monomial coefficients in \(n\) are all integers is needed.
Also, the \(C_{k,j}\) generally have alternating signs in \(t\); positivity
in the conjecture concerns \(R_{k,n}\), not these numerator blocks.
\end{remark}

\begin{corollary}[Resolution of the selected conjecture]
Conjecture~8 of Cigler's paper \cite{cigler} holds for every
\(k\ge1,n\ge0\). Its polynomiality, positive full support, exact degree,
palindromicity, unimodality, numerator representation, numerator degrees,
and numerator reciprocity are established respectively by
Theorem~\ref{thm:main}, Corollary~\ref{cor:support},
Theorem~\ref{thm:unimodal}, and Theorem~\ref{thm:numerator}.
\end{corollary}

\section{Coefficient edges and special values}
\label{sec:consequences}

\begin{proposition}[Stable initial coefficients]
\label{prop:edge}
For \(k\ge1\) and \(0\le r\le n\),
\begin{equation}
 [t^r]R_{k,n}(t)=\binom{k^2+r-1}{r}.
 \label{eq:edge}
\end{equation}
The next coefficient, with coefficients outside the degree interpreted as
zero, is
\begin{equation}
 [t^{n+1}]R_{k,n}(t)
 =\binom{k^2+n}{n+1}-\binom{n+k}{k-1}^{\!2}.
 \label{eq:edge-next}
\end{equation}
\end{proposition}
\begin{proof}
The \(j=0\) term in \eqref{eq:numerator} is one; every other term starts
at degree at least \(n+1\). Thus
\[
 R_{k,n}(t)\equiv(1-t)^{-k^2}\pmod{t^{n+1}},
\]
which proves \eqref{eq:edge} by the binomial series.
At degree \(n+1\), only the constant coefficient of \(C_{k,1}\) supplies
a correction. Its unique subset has lower indices \(0,\ldots,k-2\) and
upper index \(k\). Evaluation of \eqref{eq:weight} gives
\[
 C_{k,1}(n,0)
 =\left(\frac{(n+2)(n+3)\cdots(n+k)}{(k-1)!}\right)^2
 =\binom{n+k}{k-1}^{\!2}.
\]
Empty products cover \(k=1\). Substitution in \eqref{eq:numerator} proves
\eqref{eq:edge-next}.
\end{proof}

\begin{proposition}[Dimension product and monotonicity in the order]
\label{prop:dimension}
For \(k,n\ge0\),
\begin{equation}
 R_{k,n}(1)=\prod_{i=1}^{k}\prod_{j=1}^{k}
 \frac{n+i+j-1}{i+j-1}.
 \label{eq:dimension}
\end{equation}
Also, \(R_{k,n+1}(t)-R_{k,n}(t)\) has nonnegative coefficients, and
\(R_{k,n}(0)=1\).
\end{proposition}
\begin{proof}
At \(t=1\), Theorem~\ref{thm:main} gives
\(s_{(n^k)}(1^{2k})\). In the dimension formula
\eqref{eq:weyl-dimension}, now with \(2k\) letters, factors from pairs
within the same half cancel. The remaining factors are
\[
 \prod_{a=1}^{k}\prod_{b=k+1}^{2k}\frac{n+b-a}{b-a}.
\]
Set \(i=k+1-a\) and \(j=b-k\) to obtain \eqref{eq:dimension}.
The rectangles \((n^k)\) form an increasing family, and the summand
\(d_k(\lambda)^2t^{|\lambda|}\) in \eqref{eq:positive} is independent
of \(n\). This proves coefficientwise monotonicity. The last assertion
was already proved in Corollary~\ref{cor:support}.
\end{proof}

The apparent poles at \(t=1\) in \eqref{eq:numerator} are therefore
removable. For the original determinant,
\(D_{2k}(n,1)=R_{k,n}(1)\). At \(t=0\), the original determinant is zero
for \(n\ge2\), while its normalized polynomial has value one. These are
different statements and should not be conflated.

The parameter \(t\) records how many tableau entries are in the upper
half of the alphabet. This is a repeated-variable specialization, not the
principal specialization \(1,t,t^2,\ldots\); standard volume-generating
formulas for plane partitions cannot simply be substituted for it.

\section{A generic recurrence and three asymptotic regimes}
\label{sec:recurrence}
The explicit numerator has information beyond coefficient shape. The
following results concern fixed \(k\) and varying matrix order \(n\).

\begin{theorem}[Exact generic recurrence]
\label{thm:recurrence}
Fix \(k\ge1\), and regard \(t\) as an indeterminate. Let \(E\) be the
forward shift in \(n\). Over \(\Q(t)\), the minimal monic
constant-coefficient recurrence operator for \(R_{k,n}(t)\) is
\begin{equation}
 \prod_{j=0}^{k}(E-t^j)^{2j(k-j)+1}.
 \label{eq:recurrence}
\end{equation}
Equivalently, the generating function in the order \(n\) has reduced
denominator
\begin{equation}
 \prod_{j=0}^{k}(1-t^jz)^{2j(k-j)+1}.
 \label{eq:gf-denominator}
\end{equation}
Its degree in \(z\), and the recurrence order, is
\begin{equation}
 L_k=\sum_{j=0}^{k}\bigl(2j(k-j)+1\bigr)
     =\frac{(k+1)(k^2-k+3)}{3}.
 \label{eq:recurrence-order}
\end{equation}
\end{theorem}
\begin{proof}
Rewrite \eqref{eq:numerator} in the form
\[
 R_{k,n}(t)=\sum_{j=0}^{k}t^{jn}Q_j(n,t),\qquad
 Q_j(n,t)=\frac{(-1)^jt^{j^2}C_{k,j}(n,t)}{(1-t)^{k^2}}.
\]
By Theorem~\ref{thm:numerator}, \(Q_j\) has exact degree
\(d_j=2j(k-j)\) in \(n\). Finite differences give
\((E-t^j)^{d_j+1}(t^{jn}Q_j(n,t))=0\); multiplying all these commuting
operators proves the recurrence.

For minimality, recall that the generating function of a polynomial in
\(n\) of exact degree \(d\) has a pole of exact order \(d+1\) at one.
This follows by applying \((w\,d/dw)^d\) to \((1-w)^{-1}\): the
coefficient of its highest-order pole is \(d!\), times the leading
coefficient of the polynomial. The \(j\)-th summand above consequently
has a pole of exact order \(d_j+1\) at \(z=t^{-j}\). Over \(\Q(t)\)
these locations are distinct, so no other summand cancels that pole.
This proves the exact denominator and hence minimality. Summing the
exponents gives \eqref{eq:recurrence-order}.
\end{proof}

The denominator is the one appearing in the nearby Conjecture~10 of
\cite{cigler}. That conjecture also specifies a numerator degree and a
separate reciprocity identity, which are not asserted here. The present
article's complete-resolution claim concerns Conjecture~8 only.
For specialized values of \(t\), poles can coincide or disappear, so
Theorem~\ref{thm:recurrence} must not be read as an unchanged minimal-order
claim after every specialization. At \(t=0\), the sequence is constant;
at \(t=1\), it is a polynomial in \(n\) of degree \(k^2\).

\Needspace{17\baselineskip}
\begin{theorem}[Fixed-parameter asymptotics]
\label{thm:asymptotics}
Fix \(k\ge1\). For every fixed real \(t\) with \(0<t<1\),
\begin{equation}
 R_{k,n}(t)\longrightarrow(1-t)^{-k^2}
 \quad(n\to\infty),
 \label{eq:limit}
\end{equation}
and the deficit has the sharp leading term
\begin{equation}
 (1-t)^{-k^2}-R_{k,n}(t)
 \sim\frac{n^{2k-2}t^{n+1}}
 {((k-1)!)^2(1-t)^{k^2-2k+2}}.
 \label{eq:deficit}
\end{equation}
At \(t=1\),
\begin{equation}
 R_{k,n}(1)\sim
 \frac{n^{k^2}}{\displaystyle\prod_{i=1}^{k}\prod_{j=1}^{k}(i+j-1)}.
 \label{eq:asymp-one}
\end{equation}
For fixed \(t>1\),
\begin{equation}
 R_{k,n}(t)\sim t^{kn}(1-t^{-1})^{-k^2}.
 \label{eq:asymp-large}
\end{equation}
\end{theorem}
\begin{proof}
For \(0<t<1\), every term with \(j\ge1\) in the numerator formula
vanishes as \(n\to\infty\), because it is a fixed-degree polynomial in
\(n\) times \(t^{jn}\). This proves \eqref{eq:limit}.

To identify the first error term, compute the leading coefficient in
\(n\) of \(C_{k,1}\). Put \(m=k-1\). A selected subset with one
upper-half index consists of all lower-half indices except some \(a\),
and the single upper index \(k+b\), where \(0\le a,b\le m\). Its
\(\delta\) is \(m-a+b\). Deleting one element from a consecutive
Vandermonde gives
\[
 [n^{2m}]W_S(n)=\frac{\binom ma\binom mb}{(m!)^2}.
\]
Consequently,
\begin{align*}
 [n^{2k-2}]C_{k,1}(n,t)
 &=\frac{1}{(m!)^2}
   \sum_{a,b=0}^{m}\binom ma\binom mb(-t)^{m-a+b}\\
 &=\frac{(1-t)^{2m}}{(m!)^2}.
\end{align*}
The \(j=1\) term in \eqref{eq:numerator} is
\(-t^{n+1}C_{k,1}\). All terms with \(j\ge2\) are exponentially
smaller than this term for fixed \(t\in(0,1)\). Division by
\((1-t)^{k^2}\) proves \eqref{eq:deficit}. For \(k=1\), the error
formula is exact: it is \(t^{n+1}/(1-t)\).

Equation~\eqref{eq:asymp-one} follows from the product in
\eqref{eq:dimension}. Finally, palindromicity reduces the case \(t>1\)
to \eqref{eq:limit} with parameter \(t^{-1}\).
\end{proof}
These are fixed-\(t\) assertions, not estimates uniform as \(t\to1\).
A transition regime with \(t\) depending on \(n\) would require a
separate analysis.

\section{Worked examples and explicit checks on conventions}
\label{sec:examples}
For \(k=1\), the rectangle has one row. Either the positive formula or
Jacobi--Trudi gives
\begin{equation}
 R_{1,n}(t)=1+t+\cdots+t^n=\frac{1-t^{n+1}}{1-t}.
 \label{eq:k1}
\end{equation}
Here \(C_{1,0}=C_{1,1}=1\), and the numerator formula is immediate.
This family also shows why ``unimodal'' cannot be strengthened to
``strictly unimodal'' without exceptions. For \(n=2\), the roots of
\(1+t+t^2\) are nonreal; real-rootedness is not a premise of our proof.

For \(k=2\), evaluation of \eqref{eq:blocks} gives
\begin{align}
 C_{2,0}=C_{2,2}&=1,\notag\\
 C_{2,1}(n,t)&=(n+2)^2-2(n^2+4n+3)t+(n+2)^2t^2.
 \label{eq:k2C}
\end{align}
Thus
\begin{equation}
 R_{2,n}(t)=
 \frac{1-t^{n+1}\bigl((n+2)^2-2(n^2+4n+3)t+(n+2)^2t^2\bigr)
       +t^{2n+4}}{(1-t)^4}.
 \label{eq:k2}
\end{equation}
Cigler had already obtained this small-shift formula in his equation (53);
it is included as a normalization check, not claimed as a new special case.
The contribution here is the all-\(k\) mechanism and its consequences.

Some small normalized polynomials are:
\begin{center}
\small
\begin{tabular}{@{}ccl@{}}
\toprule
\(k\)&\(n\)&\(R_{k,n}(t)\)\\
\midrule
1&2&\(1+t+t^2\)\\
2&1&\(1+4t+t^2\)\\
2&2&\(1+4t+10t^2+4t^3+t^4\)\\
2&3&\(1+4t+10t^2+20t^3+10t^4+4t^5+t^6\)\\
3&1&\(1+9t+9t^2+t^3\)\\
3&2&\(1+9t+45t^2+65t^3+45t^4+9t^5+t^6\)\\
\bottomrule
\end{tabular}
\end{center}
For \(k=3,n=3\), the coefficient list is
\[
 (1,9,45,165,270,270,165,45,9,1).
\]
Its weight-string multiplicities up to the center are
\((1,8,36,120,105)\), the consecutive differences prescribed by
\eqref{eq:multiplicities}.

As a direct two-by-two check on the Hankel normalization,
\begin{align*}
 D_2(2,t)&=(1+t)(1+4t+t^2)-(1+2t)^2\\
        &=t+t^2+t^3=t^{\lfloor 2^2/4\rfloor}R_{1,2}(t).
\end{align*}
The factor \(t\) is essential; omitting it would make both the Schur
identity and the stated coefficient degree wrong.

\section{Verification, dependencies, and limits of the claim}
\label{sec:verification}

\subsection{What the executable verifier does}
The accompanying file \texttt{code/verify.py} uses only Python's standard
library and arbitrary-precision integers. Polynomials are stored as
coefficient tuples in ascending degree. It implements addition,
multiplication, and exact polynomial division, with explicit checks for
zero remainders and integral quotients. The determinant engine is the
fraction-free Bareiss algorithm with row pivoting; every intermediate
exact division is checked.

For all \(0\le k\le6\) and \(0\le n\le10\), it compares the following
formulas:
\begin{enumerate}[label=(\arabic*)]
\item the original shifted Hankel determinant divided exactly by
      \(t^{\lfloor n^2/4\rfloor}\);
\item the elementary-symmetric Toeplitz determinant
      \(\det(e_{k+i-j}(X_k))\);
\item the complete-symmetric Jacobi--Trudi determinant
      \(\det(h_{n-i+j}(X_k))_{0\le i,j<k}\);
\item the positive rectangular partition sum \eqref{eq:positive};
\item for \(k\ge1\), the subset numerator \eqref{eq:blocks} reconstructed
      through exact division by \((1-t)^{k^2}\).
\end{enumerate}
The three determinant formulations share the Bareiss implementation;
they are independent formulas, not three independent software systems.
As a separate audit, every original Hankel determinant of order at most
four in the tested grid is also calculated by the signed-permutation
Leibniz formula, which does not use Bareiss elimination.

The verifier also checks degree, positivity of every coefficient,
palindromicity, weak unimodality, the dimension product at \(t=1\), the
stable leading coefficients, coefficientwise monotonicity in \(n\), and
the exact degrees and reciprocity of the numerator blocks. The saved run
reported:
\begin{center}
\begin{tabular}{@{}lr@{}}
\toprule
Check category & Count\\
\midrule
Parameter pairs \((k,n)\) &77\\
Determinant-to-partition comparisons &231\\
Numerator reconstructions &66\\
Independent signed-permutation audits &35\\
Disagreements &0\\
\bottomrule
\end{tabular}
\end{center}
The run used Python~3.13.5. Exact outputs are saved in
\texttt{results/verification.json}, \texttt{results/polynomials.json}, and
\texttt{results/numerator\_blocks.json}. To reproduce the default grid,
run from the package directory:
\begin{verbatim}
python3 code/verify.py
\end{verbatim}
Larger bounds can be selected with \texttt{--max-k} and
\texttt{--max-n}. Enumeration of partitions is finite but grows with the
rectangle; the default run is deliberately a check rather than a large-scale
search. No floating-point computation, numerical root test, or guessed
recurrence is used to establish an identity.

The supplementary script \texttt{code/verify\_consequences.py} adds
66 exact checks of \eqref{eq:edge-next}, six symbolic checks of the
leading coefficient controlling \eqref{eq:deficit}, and nine recurrence
windows for \(k=1,2,3\), computed from original Hankel determinants up
to order 14. All passed. The leading coefficients are checked by exact
finite differences in \(n\), not by numerical fitting. The output is
\texttt{results/consequences.json}; reproduce it with
\begin{verbatim}
python3 code/verify_consequences.py
\end{verbatim}

\subsection{Proof audit}
The roles of the main ingredients are distinct. The moment identity and
basis change prove the full Hankel-to-Toeplitz reduction, including the
normalization. Dual Jacobi--Trudi identifies the resulting determinant.
The tableau decomposition proves full support and palindromicity without
representation theory. The Schur-module character theorem and the proved
weight-string lemma give unimodality. The confluent alternant and its
Laplace expansion prove the numerator assertions, including their signs
and exact degrees. The generic recurrence and asymptotics are consequences
of those explicit formulas.

In particular, no induction step divides by a shifted Hankel determinant
that might vanish at a special parameter value. Intermediate calculations
in \(\Q(t)\) are converted to polynomial identities before specialization.
The proof treats both parities of \(n\), the empty determinant, \(k=0\),
\(n=0\), and the apparent singularities at \(t=0\) and \(t=1\).

The imported mathematical facts are the tableau/Schur identities
\eqref{eq:jt}--\eqref{eq:alternant} and the classical Schur-module
character theorem. These are standard theorems with references, not
unproved conjectures or results borrowed from the supplied manifest.
The finite checks are supplemental and cannot replace the all-parameter
proof.

\subsection{What is not asserted}
This article does not settle the odd-shift Conjecture~9, the remaining
numerator assertions of Conjecture~10, or Conjectures~13--15 concerning
the separate family catalogued in the manifest. It makes no general
real-rootedness or strict-unimodality assertion. It does not claim that the
small examples already printed by Cigler are new.

The proof has been derived and checked within this research session, but
there has been no external referee report and no proof-assistant
formalization. The source version explicitly formulates the selected
problem as a conjecture, and the search did not locate a prior resolution;
neither fact rules out an unpublished proof, a differently indexed result,
or a resolution outside the sources searched. The mathematical statement
and its argument should be assessed independently of any priority claim.

\section{Conclusion}
The even-shift middle-binomial Hankel determinant conceals a particularly
simple Toeplitz determinant. The parity-sensitive map
\(x^{2a}\mapsto A^a\), \(x^{2a+1}\mapsto(1+z)A^a\), together with the
involution \(z\mapsto t/z\), exposes it. Once the precise power
\(t^{\lfloor n^2/4\rfloor}\) is accounted for, the normalized answer is a
rectangular Schur polynomial at two equally repeated values.

That one identification separates the conjecture into three transparent
mechanisms: tableaux for positive symmetric coefficients, \(\SL_2\)
strings for unimodality, and complementary minors of a confluent alternant
for the rational numerator and its reciprocity. Each part works for all
\(k\) and \(n\); none is an extrapolation from the finite test grid.
This supplies a complete proof draft of the selected, nonduplicate
Conjecture~8, together with explicit formulas and reproducible checks.

\begin{thebibliography}{9}
\bibitem{cigler}
Johann Cigler,
\emph{Hankel determinants of middle binomial coefficients and conjectures for
some polynomial extensions and modifications},
arXiv:2111.14492v3, 30 December 2021.
Section~4, Conjecture~8, printed pp.~16--17; equations (55)--(57).
\url{https://arxiv.org/abs/2111.14492v3}.
Accessed 20 September 2026.

\bibitem{fulton}
William Fulton,
\emph{Young Tableaux: With Applications to Representation Theory and Geometry},
London Mathematical Society Student Texts~35, Cambridge University Press.
Parts I--II; in particular the Schur-module character statement in the
introduction to Part~II, pp.~79--82.
Online chapter DOI:
\url{https://doi.org/10.1017/CBO9780511626241.010}.
Publisher's character-theorem summary accessed 20 September 2026.

\bibitem{macdonald}
I.~G. Macdonald,
\emph{Symmetric Functions and Hall Polynomials}, second edition,
Oxford University Press, 1995.
Chapter~I, especially the Schur-function identities and specializations.

\bibitem{manifest}
Vladimir Reshetnikov,
\emph{A manifest of \texttt{docs/reports}}, 20 September 2026.
User-supplied \LaTeX{} file \texttt{manifest.tex}, catalogue of 71 research
packages. The related entry is
``Product formulas for ballot-polynomial Hankel determinants'', concerning
Cigler's Conjectures~13--15. Used for topic selection and duplicate avoidance,
not as a proof source.
\end{thebibliography}
\end{document}
